% LaTeX Canvas: Chapter Outline
\chapter{Jupyter}
\label{ch:jupyter}

\section*{Purpose}
Jupyter notebooks are an effective medium for exploratory analysis because they combine code, results, and narrative. The same flexibility can also create confusion: wrong working directories, out-of-order execution, hidden state, and notebooks that cannot be reproduced. This chapter teaches novices how to launch Jupyter correctly, navigate files, use cells effectively, run shell commands safely, and adopt notebook discipline so work remains interpretable and reproducible.

\section*{Learning objectives}
By the end of this chapter, you should be able to:
\begin{enumerate}
\item Launch Jupyter Notebook or JupyterLab from the command line in the correct folder.
\item Diagnose the common ``Jupyter has no files'' problem (wrong working directory).
\item Explain kernels, sessions, and the difference between Notebook and JupyterLab.
\item Use cell types (code/markdown/raw), move cells, and manage execution order.
\item Use IPython magics and run shell commands inside a notebook safely.
\item Structure a notebook for readability: headings, narrative, and controlled outputs.
\item Make notebooks reproducible: restart-and-run-all, environment capture, and data provenance.
\item Convert notebooks into scripts/reports when appropriate.
\end{enumerate}

\section*{Running theme: notebooks are documents \emph{and} programs}
A notebook should read like a report and execute like a program. The discipline is to ensure both.

\section{Mental models: what Jupyter is doing}
\subsection{Notebook vs JupyterLab}
\begin{itemize}
\item \textbf{Jupyter Notebook} (classic) is a single-document interface.
\item \textbf{JupyterLab} is a multi-document environment (files, terminals, notebooks, consoles).
\item Both run on the same concepts: a server, a browser UI, and kernels.
\end{itemize}

\subsection{Server, browser, kernel}
\begin{itemize}
\item \textbf{Server:} the process you start (the thing listening on \texttt{localhost:8888}).
\item \textbf{Browser UI:} where you edit and run cells.
\item \textbf{Kernel:} the language runtime executing your code (Python, R, etc.).
\end{itemize}

\subsection{State and order}
\begin{itemize}
\item Notebooks allow executing cells out of order.
\item The kernel remembers variables from prior cells.
\item ``It works on my notebook'' often means hidden state.
\end{itemize}

\section{Launching Jupyter the right way}
\subsection{The rule: start from your project folder}
\begin{itemize}
\item Open a terminal.
\item \textbf{Navigate} to the folder containing your notebooks (your project root).
\item Start Jupyter there.
\end{itemize}

\subsection{Launch commands (conceptual)}
\begin{itemize}
\item \texttt{jupyter lab}
\item \texttt{jupyter notebook}
\item Common pattern: pass a directory path if you are not already in it.
\end{itemize}

\subsection{Verifying you launched in the correct directory}
\begin{itemize}
\item In the terminal: confirm the current directory before launching.
\item In the UI: confirm the file browser shows your expected folder tree.
\item Create a small marker file (e.g., \texttt{PROJECT\_ROOT.txt}) so you can tell when you are in the right place.
\end{itemize}

\subsection{Launching from an environment (important for students)}
\begin{itemize}
\item Ensure you started the correct environment (conda/venv) before launching Jupyter.
\item Know how to confirm which Python/kernel is in use (kernel name, version cell).
\end{itemize}

\section{The common failure: ``Jupyter shows no notebooks''}
\subsection{Symptom}
\begin{itemize}
\item The Jupyter file browser is empty or shows a different folder than expected.
\item Your notebook cannot find \texttt{data/} even though it exists.
\end{itemize}

\subsection{Root cause: wrong working directory / wrong server}
\begin{itemize}
\item You launched Jupyter from a different folder.
\item You have multiple Jupyter servers running and you opened the wrong tab.
\item You launched from the wrong environment, so kernels and paths differ.
\end{itemize}

\subsection{Fix checklist}
\begin{enumerate}
\item Stop Jupyter (Ctrl+C in the terminal) and close the browser tab.
\item In the terminal: \textbf{navigate} to the correct folder.
\item Relaunch Jupyter.
\item If still wrong: check for other running Jupyter servers and shut them down.
\item Add a first cell in notebooks that prints the working directory.
\end{enumerate}

\section{Notebook mechanics: cells and execution}
\subsection{Cell types}
\begin{description}
\item[Code] Executes in the kernel.
\item[Markdown] Narrative text, headings, equations, images.
\item[Raw] Passed through without formatting (rare; use intentionally).
\end{description}

\subsection{Running cells}
\begin{itemize}
\item Run one cell vs run all.
\item Interrupt vs restart kernel.
\item Understand the execution count as evidence of order.
\end{itemize}

\subsection{Moving and organizing cells}
\begin{itemize}
\item Move cells to keep the story and the computation aligned.
\item Use headings to create a navigable outline.
\item Keep imports and configuration near the top.
\end{itemize}

\subsection{Output management}
\begin{itemize}
\item Avoid printing entire datasets.
\item Prefer summaries (shape, head, descriptive stats).
\item Clear and re-run outputs as part of reproducibility checks.
\end{itemize}

\section{Running shell commands inside notebooks (responsibly)}
\subsection{Why this is useful}
\begin{itemize}
\item Quick checks: list files, inspect paths, verify data presence.
\item Lightweight automation: download a file, decompress, run a command-line tool.
\end{itemize}

\subsection{Two mechanisms: ``bang'' and magics}
\begin{itemize}
\item \textbf{Bang commands} (\texttt{!}): run a shell command from a code cell.
\item \textbf{Magics} (\texttt{%} and \texttt{%%}): special IPython commands for workflows.
\end{itemize}

\subsection{Common magics for students}
\begin{itemize}
\item Directory awareness: \texttt{%pwd}, \texttt{%cd}
\item Timing: \texttt{%time}, \texttt{%timeit}
\item Multi-line shell blocks (when taught): \texttt{%%bash} (or equivalent)
\end{itemize}

\subsection{Safety rules for shell commands in notebooks}
\begin{itemize}
\item Treat shell commands as potentially destructive.
\item Avoid \texttt{sudo} from notebooks.
\item Never embed secrets (tokens, passwords) in notebook cells.
\item Record what the command does and why (markdown cell above).
\item Prefer reproducible commands over manual clicking.
\end{itemize}

\section{Notebook style and discipline (how to write a notebook that survives)}
\subsection{Structure: a notebook template}
\begin{itemize}
\item Title + one-paragraph purpose.
\item Setup: imports, configuration, paths.
\item Data acquisition and provenance.
\item Cleaning and validation checks.
\item Analysis/modeling.
\item Results and interpretation.
\item Next steps / limitations.
\end{itemize}

\subsection{Narrative: make the notebook readable}
\begin{itemize}
\item Use markdown headings and short explanatory paragraphs.
\item Explain assumptions and decisions (why this filter? why this join?).
\item Use captions on plots and tables.
\end{itemize}

\subsection{Reproducibility discipline}
\begin{itemize}
\item Use a consistent top-to-bottom execution order.
\item Periodically: \textbf{Restart kernel and Run All}.
\item Avoid hidden dependencies on prior interactive state.
\item Prefer functions over repeated copy/paste code.
\item Prefer relative paths inside a project.
\end{itemize}

\subsection{Environment capture (student level)}
\begin{itemize}
\item Record: Python version, key package versions, and OS.
\item Store environment files (e.g., \texttt{environment.yml}) when using conda.
\item Keep datasets and generated artifacts in well-labeled folders.
\end{itemize}

\subsection{Outputs and size control}
\begin{itemize}
\item Large outputs make notebooks slow and hard to version.
\item Avoid embedding huge binary blobs.
\item Prefer saving outputs to files (\texttt{data/processed}, \texttt{figures/}).
\end{itemize}

\section{Notebook pitfalls and how to prevent them}
\subsection{Out-of-order execution}
\begin{itemize}
\item Prevention: restart-and-run-all checks.
\item Signal: execution counts jump around; variables appear ``from nowhere.''
\end{itemize}

\subsection{Wrong kernel / wrong environment}
\begin{itemize}
\item Symptom: imports fail or versions differ from expectations.
\item Prevention: name kernels clearly; confirm versions in a top cell.
\end{itemize}

\subsection{File not found (paths)}
\begin{itemize}
\item Symptom: \texttt{FileNotFoundError} despite file existing.
\item Prevention: print working directory; use project-relative paths.
\end{itemize}

\subsection{Long-running or stuck cells}
\begin{itemize}
\item Learn: interrupt vs restart.
\item Add progress indicators or timing.
\item Move expensive work to scripts or batch runs when appropriate.
\end{itemize}

\section{When to move from notebooks to scripts (and back)}
Chapter~\ref{ch:scripts-vs-notebooks} covers the scripting workflow in full; this section explains how notebooks and scripts fit together.
\subsection{Notebook strengths}
\begin{itemize}
\item Exploration, explanation, and sharing results.
\end{itemize}

\subsection{Script strengths}
\begin{itemize}
\item Production runs, automation, CI testing, and parameterized workflows.
\end{itemize}

\subsection{A practical boundary for students}
\begin{itemize}
\item Keep exploration in notebooks.
\item Extract stable logic into \texttt{src/} as functions.
\item Keep the notebook as a narrative driver calling those functions.
\end{itemize}

\section{Worked examples (outline)}
\subsection{Example 1: Launch Jupyter in the right place}
\begin{itemize}
\item Navigate to a project folder.
\item Start JupyterLab.
\item Create a notebook and verify the working directory.
\end{itemize}

\subsection{Example 2: Debug ``no files'' / wrong directory}
\begin{itemize}
\item Demonstrate launching from the wrong folder.
\item Diagnose and fix by relaunching from the correct path.
\item Add a working-directory check cell.
\end{itemize}

\subsection{Example 3: Shell commands for file checks}
\begin{itemize}
\item List project files.
\item Confirm a dataset exists and report size.
\item Use \texttt{%pwd} and \texttt{%cd} intentionally.
\end{itemize}

\subsection{Example 4: Make a notebook reproducible}
\begin{itemize}
\item Add headings and narrative.
\item Extract repeated code into functions.
\item Restart kernel and run all; fix hidden-state bugs.
\end{itemize}

\section{Templates}
\subsection{Template A: Notebook header block}
\begin{verbatim}
Title:
Author:
Date:
Purpose:
Data sources:
Environment:

* Python:
* Key packages:
  How to run:
* Restart kernel and Run All
  Outputs:
* Where figures/tables are saved
  \end{verbatim}

\subsection{Template B: Reproducibility checklist cell}
\begin{verbatim}

# Reproducibility check

# 1) print working directory

# 2) print python and package versions

# 3) confirm data files exist

# 4) run a small smoke test

\end{verbatim}

\section{Exercises}
\begin{enumerate}
\item Launch JupyterLab from a project directory and create a notebook that prints its working directory.
\item Intentionally launch from the wrong directory; diagnose and fix.
\item Practice cell operations: convert code to markdown, move cells, and create headings.
\item Use one shell command (via \texttt{!}) to list files and one magic (\texttt{%pwd}) to confirm location.
\item Take a messy notebook and refactor it into sections with narrative and a restart-and-run-all proof.
\item Extract one repeated block into a function and verify outputs are unchanged.
\end{enumerate}

\section{One-page checklist}
\begin{itemize}
\item I launch Jupyter from the correct project folder (or pass the correct directory).
\item I can diagnose ``no files'' as a working-directory/server mix-up.
\item I understand kernels and can choose the correct one.
\item I use markdown headings to structure the notebook.
\item I keep a top-to-bottom execution order and periodically restart-and-run-all.
\item I use shell commands and magics sparingly, with explanations and no secrets.
\item I record environment details and keep data provenance explicit.
\item I keep outputs controlled and store large artifacts outside the notebook.
\end{itemize}

% \appendix
\section{Quick reference: common launch and debugging moves}
\begin{itemize}
\item Confirm working directory before launching.
\item If you see the wrong files: stop server, relaunch in correct directory.
\item If imports fail: check kernel/environment.
\item If execution hangs: interrupt; then restart if needed.
\end{itemize}

\section{Quick reference: IPython conveniences}
\begin{itemize}
\item \texttt{!command} runs a shell command.
\item \texttt{%pwd} shows current directory; \texttt{%cd} changes it.
\item Use timing magics to understand slow cells.
\end{itemize}
